\section*{الفصل \ref{ch:b2:biomolecules} --- الجزيئات الحيوية: الأحماض الأمينية والسكريات والليبيدات والنوكليوتيدات}

\begin{solution}{exo:b2:biomolecules:1}
\begin{itemize}
  \item بين pH 3 و9، \omterm{def:b2:biomolecules:amino-acid}{الأيون المزدوج الشحنة} \ce{H3N+-CH(CH3)-COO-}.
  \item عند pH 1، الكاتيون \ce{H3N+-CH(CH3)-COOH}.
  \item عند pH 12، الأنيون \ce{H2N-CH(CH3)-COO-}.
\end{itemize}
\end{solution}

\begin{solution}{exo:b2:biomolecules:2}
الغليسين: $(2.35 + 9.78)/2 \approx 6.1$. الألانين: $(2.35 + 9.87)/2 \approx 6.1$.
% ledger: pka:glycine1, pka:glycine2, pka:alanine1, pka:alanine2
\end{solution}

\begin{solution}{exo:b2:biomolecules:3}
\ce{H2N-CH(CH3)-CO-NH-CH2-COOH}: يقدم الألانين الطرف N، والغليسين الطرف C. أما Gly--Ala،
\ce{H2N-CH2-CO-NH-CH(CH3)-COOH}، فمتماكب بنيوي: المجموعة الأمينية الحرة فيه على الغليسين.
\end{solution}

\begin{solution}{exo:b2:biomolecules:4}
\ce{NH2} إلى اليسار: L. الأولويات \babelsublr{\ce{NH2} > \ce{COOH} > \ce{CH2OH} > \ce{H}}: فالسيرين L
هو (\textit{S}).
\end{solution}

\begin{solution}{exo:b2:biomolecules:5}
المجموعات: أمونيوم الطرف N (نحو 9)، \omterm{def:b2:biomolecules:amino-acid}{والسلسلة الجانبية} لليسين (10.7)، \omterm{def:b2:biomolecules:amino-acid}{والسلسلة الجانبية} للأسبارتات (3.9)،
والمجموعة الكربوكسيلية للطرف C (نحو 2). عند pH 1: $+1 +1 + 0 + 0 = +2$. وعند pH 7: $+1 +1 -1 -1 = 0$. وعند pH 12:
$0 + 0 - 1 - 1 = -2$ (\omterm{def:b2:biomolecules:amino-acid}{فالسلسلة الجانبية} لليسين، الواقعة 1.3 وحدة تحت pH 12، متعادلة في معظمها).
\end{solution}

\begin{solution}{exo:b2:biomolecules:6}
كلوريد الثيونيل والتسخين كانا سيهاجمان مجموعات أخرى أيضًا، والأهم أن \omterm{def:b2:acyl-substitution:derivative}{كلوريد الأسيل} الشديد التفاعلية
للحمض الأميني يفقد ترتيبه الفراغي بسهولة (إذ يصير هيدروجين $\alpha$ حمضيًا): فيصير \omterm{def:b2:biomolecules:peptide}{الببتيد}
راسيميًا جزئيًا. أما \omterm{def:b2:biomolecules:peptide}{عامل الاقتران} فينشط الحمض في درجة حرارة الغرفة، في ظروف لطيفة.
\end{solution}

\begin{solution}{exo:b2:biomolecules:7}
الذرة O5 مرتبطة بالذرة C2: حلقة من C2 وC3 وC4 وC5 وO، أي خمس ذرات (فورانوز). وبإسقاط هاورث، يكون O الحلقة
في الخلف، وC2 على اليمين مع \ce{CH2OH} (C1) إلى الأسفل و\ce{OH} إلى الأعلى ($\beta$)؛ وعلى C3 \ce{OH} إلى الأعلى (يسارًا في
فيشر)، وعلى C4 \ce{OH} إلى الأسفل، وC5 يحمل \ce{CH2OH} (C6) إلى الأعلى.
\end{solution}

\begin{solution}{exo:b2:biomolecules:8}
$x_\alpha = (80.2 - 52.8)/(150.7 - 52.8) = 27.4/97.9 \approx 0.28$: 28~\% من $\alpha$، و72~\% من $\beta$.
\end{solution}

\begin{solution}{exo:b2:biomolecules:9}
يحتفظ المالتوز بنصف أسيتال حر على غلوكوزه الثاني، ينفتح إلى ألدهيد: فهو مختزل. وفي السكروز
ينخرط الكربونان الأنوميريان كلاهما في \omterm{def:b2:biomolecules:glycoside}{الرابطة الغليكوزيدية}: فلا نصف أسيتال، ولا اختزال. والحمض المخفف
يحلمه \omterm{def:b2:biomolecules:glycoside}{الرابطة الغليكوزيدية}، وهي أسيتال: فيتحرر الغلوكوز والفركتوز، وكلاهما يختزل.
\end{solution}

\begin{solution}{exo:b2:biomolecules:10}
التريولين، \qty{885.45}{g/mol}، يستهلك ثلاثة \ce{KOH} (\qty{56.105}{g/mol}): $3 \times 56.105/885.45 =
0.190$~g لكل g، أي قرينة \omterm{def:b2:acyl-substitution:saponification}{تصبّن} قدرها 190. والسلاسل الأقصر تعني كتلة مولية أصغر مع مجموعات
الإستر الثلاث نفسها: فمولات أكثر من الدهن لكل غرام، وKOH أكثر.
% ledger: aw:C, aw:H, aw:O, aw:K
\end{solution}

\begin{solution}{exo:b2:biomolecules:11}
احمِ أمين الغليسين على هيئة كربامات \textit{tert}-بيوتوكسي كربونيل، وحمض الألانين على هيئة
إستر ميثيلي. ثم اقرن بواسطة كربوديإيميد. ثم أزل الكربامات بالحمض والإستر بقاعدة مخففة
(ظروف متعامدة، فيمكن تحرير أي من الطرفين وحده): Gly--Ala.
\end{solution}

\begin{solution}{exo:b2:biomolecules:12}
\babelsublr{5$'$-ACGCAT-3$'$} (يُقرأ عكسيًا: A يتزاوج مع T، وG مع C). الأزواج: A--T، T--A، G--C، C--G، G--C، T--A:
$3 \times 2 + 3 \times 3 = 15$ رابطة هيدروجينية.
\end{solution}

\begin{solution}{pb:b2:biomolecules:1}
\textbf{1.} حمض الأسبارتيك L، والفينيل ألانين L والميثانول.
\textbf{2.} في كليهما، \ce{COOH} في الأعلى، \omterm{def:b2:biomolecules:amino-acid}{والسلسلة الجانبية} في الأسفل (\ce{CH2COOH}، \ce{CH2C6H5})، و\ce{NH2}
إلى اليسار.
\textbf{3.} (\textit{S}) لكليهما.
\textbf{4.} مركزان فراغيان: أربعة متماكبات فراغية.
\textbf{5.} \ce{H2N-CH(CH2COOH)-CO-NH-CH(CH2C6H5)-COOCH3}: \omterm{def:b2:biomolecules:peptide}{الرابطة الببتيدية} تصل
البقيتين؛ والإستر هو \ce{COOCH3} عند الطرف C.
\textbf{6.} مستقبلات الطعم كيرالية: فهي ترتبط بمتماكب فراغي واحد لا بصورته المرآوية ولا
بمتماكباته اللامرآوية، كما قالت المقدمة.
\textbf{7.} \ce{NH3+} للطرف N و\ce{COOH} في \omterm{def:b2:biomolecules:amino-acid}{السلسلة الجانبية} لحمض الأسبارتيك؛ أما حمض الطرف C
فإستر ميثيلي، بلا هيدروجين حمضي.
\textbf{8.} عند pH 2: \ce{NH3+} ($+1$)، و\ce{COOH} \omterm{def:b2:biomolecules:amino-acid}{السلسلة الجانبية} (0): $+1$.
\textbf{9.} عند pH 7: $+1 - 1 = 0$.
\textbf{10.} عند pH 12: $0 - 1 = -1$.
\textbf{11.} بين قيمتي $\mathrm pK_a$ للمجموعتين اللتين تتغيران: $(3.9 + 9.9)/2 \approx 6.9$.
\textbf{12.} في ثنائي \omterm{def:b2:biomolecules:peptide}{الببتيد} يحمل كربون الأمين أميدًا بدلًا من كربوكسيلات: فتتغير
المجموعات والشحنات المجاورة، ومعها قيم $\mathrm pK_a$ (فأمونيوم الطرف N في
\omterm{def:b2:biomolecules:peptide}{الببتيد} أكثر حموضة بوضوح من أمونيوم الحمض الأميني الحر). والنموذج يعطي الشكل، لا
القيم المضبوطة.
\textbf{13.} الحمض والأمين يكوّنان ملحًا أولًا؛ والتسخين كان سيعطي خلائط (Asp--Asp، وسلاسل أطول،
والمجموعة الكربوكسيلية الخطأ) ويحوّل المراكز الفراغية إلى راسيمية.
\textbf{14.} مجموعته الأمينية ومجموعة \omterm{def:b2:biomolecules:amino-acid}{سلسلته الجانبية} الكربوكسيلية.
\textbf{15.} يحوّل \ce{OH} الحمض إلى مجموعة مغادرة جيدة، أي مشتق منشط يهاجمه
الأمين في درجة حرارة الغرفة (استبدال أسيلي).
\textbf{16.} ثنائي \omterm{def:b2:biomolecules:peptide}{الببتيد} $\beta$، الموصول عبر \omterm{def:b2:biomolecules:amino-acid}{السلسلة الجانبية}: متماكب للأسبارتام، لا الأسبارتام.
\textbf{17.} الأمين على هيئة كربامات بنزيل أوكسي كربونيل \omterm{def:b2:biomolecules:amino-acid}{والسلسلة الجانبية} على هيئة إستر بنزيلي: يُزالان
كلاهما معًا بالتحلل الهيدروجيني على البلاديوم، الذي يترك الإستر الميثيلي سليمًا (فهما متعامدان
معه).
\textbf{18.} للحمض المنشط هيدروجين $\alpha$ حمضي؛ تنزعه القاعدة فيصير المركز الفراغي
راسيميًا.
\textbf{19.} احمِ حمض الأسبارتيك (كربامات على N، وإستر بنزيلي على \omterm{def:b2:biomolecules:amino-acid}{السلسلة الجانبية})؛ ثم نشّط
\ce{COOH} الحرة في الموضع $\alpha$ \omterm{def:b2:biomolecules:peptide}{بعامل الاقتران}؛ ثم أضف الإستر الميثيلي للفينيل ألانين L؛ ثم أجرِ التحلل الهيدروجيني.
\textbf{20.} الإستر الميثيلي \omterm{def:b2:biomolecules:peptide}{والرابطة الببتيدية}؛ والإستر أولًا، لأن الأميد أقل تفاعلية بكثير.
\textbf{21.} ثنائي \omterm{def:b2:biomolecules:peptide}{الببتيد} Asp--Phe (الحمض الحر) والميثانول.
\textbf{22.} الجزيء الحلو يُستهلك: ونواتج حلمهته لا تحل محله.
\textbf{23.} $14 \times 12.011 + 18 \times 1.008 + 2 \times 14.007 + 5 \times 15.999 =
\qty{294.307}{g/mol}$.
\textbf{24.} $0.180/294.307 = \qty{6.116e-4}{mol}$، أي \qty{0.612}{mmol}.
\textbf{25.} جزيء ميثانول واحد لكل أسبارتام: $0.6116 \times 32.042 = \boldsymbol{\approx \qty{19.6}{mg}}$ من
الميثانول.
\end{solution}
